\subsection{Imagen inversa de un $\mathcal{B}$-Módulo}
\label{subsection:0.4.3-fr}

\begin{env}[4.3.1]
\label{0.4.3.1-fr}
Con las hipótesis y la notación de (4.2.1), sean $\mathcal{G}$ un
$\mathcal{B}$-Módulo y $\psi^*(\mathcal{G})$ su imagen inversa (3.7.1), que
es, por tanto, un haz de grupos abelianos sobre $X$. La definición de las
secciones de $\psi^*(\mathcal{G})$ y de $\psi^*(\mathcal{B})$ (3.7.1)
muestra que $\psi^*(\mathcal{G})$ está provisto canónicamente de una
estructura de $\psi^*(\mathcal{B})$-Módulo. Por otra parte, el
homomorfismo $\theta^\sharp:\psi^*(\mathcal{B})\to\mathcal{A}$ dota a
$\mathcal{A}$ de una estructura de $\psi^*(\mathcal{B})$-Módulo. Para
evitar confusiones, cuando sea preciso, lo denotaremos por
$\mathcal{A}_{[\theta]}$. El producto tensorial
$\psi^*(\mathcal{G})\otimes_{\psi^*(\mathcal{B})}\mathcal{A}_{[\theta]}$
está provisto entonces de una estructura de $\mathcal{A}$-Módulo.
Diremos que este $\mathcal{A}$-Módulo es \emph{la imagen
\oldpage[0\textsubscript{I}]{41}
inversa de $\mathcal{G}$ por el morfismo}
$\Psi$ y lo denotaremos por $\Psi^*(\mathcal{G})$. Si $\mathcal{G}_1$,
$\mathcal{G}_2$ son dos $\mathcal{B}$-Módulos sobre $Y$ y $v$ es un
$\mathcal{B}$-homomorfismo $\mathcal{G}_1\to\mathcal{G}_2$,
$\psi^*(v)$, como se comprueba de inmediato, es un
$\psi^*(\mathcal{B})$-homomorfismo de $\psi^*(\mathcal{G}_1)$ en
$\psi^*(\mathcal{G}_2)$. Por consiguiente, $\psi^*(v)\otimes1$ es un
$\mathcal{A}$-homomorfismo $\Psi^*(\mathcal{G}_1)\to\Psi^*(\mathcal{G}_2)$,
que denotaremos por $\Psi^*(v)$. Así queda definido $\Psi^*$ como un
\emph{funtor covariante} de la categoría de los $\mathcal{B}$-Módulos a la
de los $\mathcal{A}$-Módulos. Aquí, a diferencia de $\psi^*$, este funtor
solo es \emph{exacto por la derecha} en general, ya que el producto tensorial
con $\mathcal{A}$ es un funtor exacto por la derecha en la categoría de los
$\psi^*(\mathcal{B})$-Módulos.

Para todo $x\in X$, se tiene
\[
  (\Psi^*(\mathcal{G}))_x
  =\mathcal{G}_{\psi(x)}\otimes_{\mathcal{B}_{\psi(x)}}\mathcal{A}_x,
\]
en virtud de (3.7.2). Por tanto, el soporte de $\Psi^*(\mathcal{G})$ está
contenido en $\psi^{-1}(\operatorname{Supp}\mathcal{G})$.
\end{env}
